% Generated draft: requires mathematical and source review.
\documentclass[11pt]{article}
\usepackage[margin=25mm]{geometry}
\usepackage{fontspec}
\setmainfont{Latin Modern Roman}
\usepackage{amsmath,amssymb,mathtools,unicode-math}
\setmathfont{Latin Modern Math}
\usepackage{enumitem,needspace,xurl,hyperref}
\setlength{\parindent}{0pt}
\setlength{\parskip}{5pt}
\setlength{\emergencystretch}{3em}
\newcommand{\N}{\mathbb{N}}
\newcommand{\Z}{\mathbb{Z}}
\newcommand{\Q}{\mathbb{Q}}
\newcommand{\R}{\mathbb{R}}
\newcommand{\C}{\mathbb{C}}
\newcommand{\F}{\mathbb{F}}
\newcommand{\A}{\mathbb{A}}
\newcommand{\PP}{\mathbb{P}}
\newcommand{\E}{\mathbb{E}}
\newcommand{\eps}{\varepsilon}
\newcommand{\dd}{\,\mathrm{d}}
\begin{document}
\section*{UnsolvedMath ID: 30006414}
\textbf{Source problem code:} catalogue entry
\subsection{Definitions and mathematical statement}
A smooth manifold is \emph{closed} if it is compact and has no boundary. A Riemannian metric \(g\) on a smooth manifold \(M\) is \emph{Einstein} if there exists a constant \(\lambda\in\R\) such that
\[
\operatorname{Ric}_g=\lambda g
\]
throughout \(M\). It has \emph{negative sectional curvature} if
\[
K_g(\sigma)<0\qquad\text{for every }p\in M\text{ and every two-dimensional subspace }\sigma\subset T_pM.
\]
For each integer \(n\geq2\), let \(\mathcal E_n\) denote the collection of closed smooth \(n\)-manifolds \(M\) for which there exist a Riemannian metric \(g\) and a constant \(\lambda\in\R\) satisfying both conditions above. In these dimensions the conditions imply \(\lambda<0\).

\textbf{Source question.} In what meaningful asymptotic or topological sense are closed manifolds admitting such metrics abundant? The collections \(\mathcal E_n\) provide notation for the question in dimensions at least two; the source does not specify a dimension range, whether dimension is fixed, or a connectedness convention.

The word ``abundant'' is unresolved in the supplied formulation. No equivalence relation on manifolds, comparison collection, counting parameter, normalization, limiting procedure, or topological criterion is specified. Consequently, the source does not determine a single quantified abundance conjecture. Choosing such data is part of making the question precise, and no particular choice is asserted here.

The source motivates the question by the difficulty of constructing Einstein metrics on closed manifolds and by a philosophy that negatively curved closed manifolds are abundant, sometimes informally described as ``most'' closed manifolds. That language supplies no defined counting measure or proportion.

The supplied literature assessment, dated 2026-08-22, labels the problem partially solved and attributes to Jaeckel and Hamenstaedt a construction of infinitely many closed manifolds with negatively curved Einstein metrics in every dimension at least four. It also states that the broader meaning of abundance is not fully classified. These are supplied curation claims, not independently verified mathematical conclusions; in particular, the supplied summary does not specify the equivalence relation under which its examples are distinct.
\subsection{Short English statement}
The question asks how plentiful compact smooth manifolds without boundary are that admit a metric which is both Einstein and strictly negatively curved on every tangent two-plane. A precise abundance claim still requires a dimension convention and a specified way to distinguish, compare, or count manifolds. The supplied dated assessment reports partial progress in dimensions at least four, but does not resolve these formulation choices.
\subsection{Sources}
\begin{itemize}
\item Supplied local UnsolvedMath registry record, permanent ID \(30006414\), source problem code \texttt{catalogue entry}, title \emph{Abundance of Negatively Curved Einstein Manifolds}, in the Oberwolfach Reports Open Problems collection. No individual proposer is identified; the recorded proposed year is 2025.
\item Supplied report citation: \emph{Differentialgeometrie im Großen (2026)}, Oberwolfach Reports, pp.\ 1591--1646, DOI \texttt{10.4171/owr/2025/30}. \url{https://doi.org/10.4171/owr/2025/30}. The literature list labels this same link \emph{Oberwolfach Report 30/2025}. The differing year indications are retained without reconciliation.
\item Supplied literature reference: \emph{Stability and construction of negatively curved Einstein metrics}. \url{https://bonndoc.ulb.uni-bonn.de/xmlui/handle/20.500.11811/13264}. The supplied assessment attributes the construction discussed above to Jaeckel and Hamenstaedt.
\item Literature assessment date: \texttt{2026-08-22}. The label ``partially solved,'' the construction claim, and the statement about the unresolved broader abundance question are supplied dated curation, not independently verified here. The record expressly requests independent verification of time-sensitive status and cited claims.
\item The original excerpt cites \texttt{[Bes08]} for background on Einstein metrics and \texttt{[So24]} in connection with the informal ``most'' philosophy. No bibliographic details for these reference labels are supplied.
\end{itemize}
\end{document}
